\documentclass{article}
\usepackage{float}
% Language setting
\usepackage[german]{babel}
\usepackage[table]{xcolor}
% Set page size and margins
% Replace `letterpaper' with `a4paper' for UK/EU standard size
\usepackage[letterpaper,top=2cm,bottom=2cm,left=3cm,right=3cm,marginparwidth=1.75cm]{geometry}
\usepackage[dvipsnames]{xcolor}
% Useful packages
\usepackage{amsmath}
\usepackage{graphicx}
\usepackage[colorlinks=true, allcolors=blue]{hyperref}
\usepackage{subcaption}
\usepackage{graphicx}
\title{Fragestellungen von Kolloquium Weisshaar I.} 
\author{Lil'Jana}
\usepackage{enumitem,amssymb}
\newlist{todolist}{itemize}{2}
\setlist[todolist]{label=$\square$}
\usepackage{pifont}
\newcommand{\cmark}{\ding{51}}%
\newcommand{\xmark}{\ding{55}}%
\newcommand{\mmark}{\ding{51}}%
\newcommand{\done}{\rlap{$\square$}{\raisebox{2pt}{\large\hspace{1pt}\cmark}}%
\hspace{-2.5pt}}
\newcommand{\maybe}{\rlap{($\square$)}{(\raisebox{2pt}{\large\hspace{1pt}\cmark}}%
\hspace{-2.5pt}}
\newcommand{\wontfix}{\rlap{$\square$}{\large\hspace{1pt}\xmark}}



\begin{document}
\maketitle
\subsection*{Die DNA, die im Gentechnik-Labor bearbeitet wird, wird meist zunächst zerkleinert durch ...}
\begin{todolist}
\item Ligasen \textit{\scriptsize Gegenteil der Restriktionsenzyme, da sie DNA-Stücke miteinander verbinden }
\item [\done] Restriktionsenzyme
\item DNA-Polymerase \textit{\scriptsize synthetisieren neue DNA-Stränge}
\item starke Säure \textit{\scriptsize sind nicht spezifisch und schädigen die DNA wahllos}
\item Röntgenstrahlen \textit{\scriptsize unkontrollierte Fragmentierung}
\end{todolist}
\subsection*{Was sind die wesentlichen Unterschiede zwischen DNA und RNA?}
\begin{itemize}
    \item Zuckerstruktur
    \item Basenpaarung
\end{itemize}
\subsection*{Welche Eigenschaften sollte ein Vektor mindestens haben?}
\begin{itemize}
    \item Replikationsursprung
    \item Selektionsmarker
    \item Klonauswahl   
\end{itemize}
\subsection*{Wofür steht die Abkürzung 'PCR'?}
Polymerase-Kettenreaktion / Polymerase Chain Reaction

\subsection*{Das fehlerhafte Protein bei der (erblichen Form der) Hypercholesterinämie ist ...}
\begin{todolist}
\item ein Enzym aus dem Aminosäure-Stoffwechsel. \textit{\scriptsize nicht direkt mit der Regulation des Cholesterinspiegels im Blut verbunden}
\item ein Enzym aus dem Lipidstoffwechsel. \textit{\scriptsize bei der erblichen Hypercholesterinämie ist nicht ein Enzym direkt betroffen, sondern der Rezeptor, } 
\item das Hämoglobin.\textit{\scriptsize hat keine direkte Rolle im Cholesterinstoffwechsel}
\item [\done] der LDL-Rezeptor. 
\item ein Enzym aus dem Kohlenhydrat-Stoffwechsel.\textit{\scriptsize auf Zuckerstoffwechsel fokussiert und nicht auf die Regulierung von Cholesterin}
\end{todolist} 
\subsection*{Wie lauten die vier fehlenden Bezeichnungen?}
\begin{figure}[H] 
    \centering
    \includegraphics[width=10cm]{bilder/nukleotid.png}
    \caption{}
    \label{fig:my_label}
\end{figure}
\subsection*{Die Funktion von Restriktionsenzymen in Bakterien ist ....}
\begin{todolist} 
\item die eigene DNA zu schneiden, um Rekombination zu ermöglichen \textit{\scriptsize erkennen normalerweise nur fremde DNA}
\item [\done] fremde DNA zu zerstören
\item fremde Proteine zu zerstören \textit{\scriptsize Restriktionsenzyme wirken auf DNA, nicht auf Proteine}
\item DNA zu synthetisieren \textit{\scriptsize Restriktionsenzyme sind nicht für die DNA-Synthese zuständig}
\item[\maybe] Plasmide zu zerschneiden \textit{\scriptsize Wenn es fremde Plasmide sind...}
\end{todolist}
\subsection*{Homologe Chromosomen unterscheiden sich häufig in den Merkmals-Informationen, die sie tragen. Sie ...}
\begin{todolist} %\textit{\scriptsize }%
\item sind deswegen unterschiedlich lang. \textit{\scriptsize Homologe Chromosomen sind in der Regel gleich lang}
\item sind deswegen beide in einer Keimzelle enthalten. \textit{\scriptsize In einer Keimzelle (Gamete), die haploid ist, ist nur eines der homologen Chromosomenpaare enthalten}
\item sind in diploiden, somatischen Zellen nicht beide vorhanden. \textit{\scriptsize In diploiden somatischen Zellen sind beide homologen Chromosomen vorhanden. Ein Chromosom jedes Paares stammt von der Mutter und eines vom Vater.}
\item [\done] enthalten deswegen öfter unterschiedliche Allele.
\end{todolist}%\textit{\scriptsize }%
\subsection*{Die tRNA, die die Aminosäure Valin translatiert, hat das Anticodon CAG.
Wie lautet das dazugehörige mRNA-Triplett?}
\begin{todolist}
\item CAG \textit{\scriptsize Anticodon der tRNA}
\item GAC \textit{\scriptsize nicht mit dem Anticodon CAG komplementär}
\item GTC \textit{\scriptsize nicht mit dem Anticodon CAG komplementär}
\item TUG \textit{\scriptsize da T (Thymin) in DNA vorkommt, aber in mRNA durch U (Uracil) }
\item [\done] GUC
\end{todolist}
\subsection*{In welchem Meiosestadium lagern sich die homologen Chromosomen zusammen, und findet "crossover" statt?}
\begin{todolist} %\textit{\scriptsize }%
\item Prophase II \textit{\scriptsize gibt keine Paarung der homologen Chromosomen und kein Crossover}
\item [\done] Prophase I
\item Metaphase II \textit{\scriptsize sind entlang der Äquatorialebene der Zelle ausgerichtet}
\item Metaphase I \textit{\scriptsize die homologen Chromosomenpaare sind entlang der Äquatorialebene der Zelle ausgerichtet, aber das Crossover findet nicht in dieser Phase statt}
\end{todolist}
\subsection*{Wie viele Typen von Keimzellen kann ein diploider Organismus mit n = 4
unabhängig vom "crossover" bilden?}
Anzahl der Chromosomenpaare: n/2 = 4/2 = 2 \newline
Mögliche Kombinationen: 2$^2$ = 4
\begin{todolist} %\textit{\scriptsize }%
\item 2
\item [\done] 4
\item 8
\item 16
\item 32
\item 64
\end{todolist}
\subsection*{Welche der genannten Aussagen ist/sind FALSCH?}
\begin{todolist} %\textit{\scriptsize }%
\item Haploide Zellen können sich durch Mitose teilen. \textit{\scriptsize richtig}
\item [\done] Haploide Zellen können sich durch Meiose teilen. \textit{\scriptsize Die Meiose ist der Prozess, bei dem diploide Zellen in haploide Zellen geteilt werden.}
\item Diploide Zellen können sich durch Mitose teilen. \textit{\scriptsize richtig}
\item Diploide Zellen können sich durch Meiose teilen. \textit{\scriptsize richtig}
\end{todolist}
\subsection*{Der Mensch hat 46 Chromosomen. Diese Anzahl (menschlicher)
Chromosomen findet man in ...}
\begin{todolist} %\textit{\scriptsize }%
\item allen Zellen des Körpers. \textit{\scriptsize nur in allen somatischen Zellen}
\item allen Eizellen und Spermien. \textit{\scriptsize Eizellen und Spermien sind haploid und enthalten 23 Chromosomen.}
\item allen Zellen der Keimbahn nach der Meiose I. \textit{\scriptsize Nach der ersten meiotischen Teilung sind die Zellen haploid und enthalten 23 Chromosomen. Diese Zellen haben die Chromosomenzahl halbiert und sind daher nicht mehr diploid. Vor der Meiose I haben die Keimbahnzellen die diploide Chromosomenanzahl von 46, aber nach der ersten meiotischen Teilung enthalten sie nur noch 23 Chromosomen.}
\item [\done] allen somatischen Zellen.
\end{todolist}
\subsection*{Ein Karyogramm wird hergestellt, in dem man ...}
\begin{todolist}%\textit{\scriptsize }%
\item den Phänotyp eines Individuums betrachtet. \textit{\scriptsize Der Phänotyp kann durch viele Faktoren beeinflusst werden, nicht nur durch die Chromosomenstruktur.}
\item den Genotyp eines Individuums analysiert. \textit{\scriptsize Karyogramm liefert es keine detaillierten Informationen über den Genotyp auf der Ebene von Einzelgenen oder deren spezifischen Allelen}
\item eine einzige Chromosomenkombination aus einer Keimzelle
betrachtet. \textit{\scriptsize Ein Karyogramm zeigt in der Regel die Chromosomen einer somatischen Zelle, nicht nur eine einzelne Keimzelle}
\item die Blutgruppe eines Individuums bestimmt. \textit{\scriptsize hat keinen direkten Zusammenhang mit dem Karyogramm}
\item[\done] die homologen Chromosomen einer Zelle nach Anzahl, Größe
und Gestalt ordnet.
\end{todolist}
\subsection*{Durch Alloploidie soll eine neue Pflanzenart erzeugt werden. Die beiden
Ausgangsspezies sind n = 4 und n = 8. Wie viele Chromosomen erwarten
sie in den somatischen Zellen der neuen Art?}
\begin{todolist}%\textit{\scriptsize }%
\item 6 
\item 12 
\item[\done] 24 \textit{\scriptsize 2×(4+8)=2×12=24}
\item 48
\item anderes
\end{todolist}
\subsection*{Was ist ein Diplont?}
\begin{todolist}%\textit{\scriptsize }%
\item[\done] Ein Organismus, in dessen Lebenszyklus es keine haploiden Zellen
gibt. \textit{\scriptsize haploide Zellen existieren nur in Form von Gameten, die zur Fortpflanzung beitragen}
\item Ein Organismus, in dessen Lebenszyklus es keine diploiden Zellen
gibt. \textit{\scriptsize Definition eines Haplonts}
\item Ein Organismus, in dessen Lebenszyklus es keine Mitose gibt. \textit{\scriptsize Mitose ist ein grundlegender Zellteilungsmechanismus}
\item Ein Organismus, in dessen Lebenszyklus es keine Meiose gibt. \textit{\scriptsize Meiose tritt in Diplonten auf}
\item[\maybe] Ein Organismus, in dessen Lebenszyklus die Produkte der Meiose in
die Befruchtung eingehen. \textit{\scriptsize trifft auf die meisten Diploiden zu, beschreibt aber nur einen Teil des Lebenszyklus }
\end{todolist}

\subsection*{Wenn Sie Crossover bzw. den Austausch von genetischem Material in
einer Tetrade zwischen Chromatiden berücksichtigen, was wird in der
Meiose I voneinander getrennt?}
\begin{todolist} %\textit{\scriptsize }%
\item[\done] Vollständige homologe Chromosomen.
\item Schwesterchromatiden von Zwei-Chromatid-Chromosomen. \textit{\scriptsize werden erst in der Meiose II getrennt}
\item Chromosomen mit gleicher Genausstattung. \textit{\scriptsize Während der Meiose I werden Chromosomen mit ähnlicher Genausstattung (homologe Chromosomen) getrennt}
\item Nur Telomere. \textit{\scriptsize  sind nicht das, was in der Meiose I getrennt wird}
\item Nur Centromere. \textit{\scriptsize Die Centromere werden erst in der Meiose II getrennt und auch dann nicht nur sie}
\item Die Stränge einer doppelsträngigen DNA in zwei Einzelstränge. \textit{\scriptsize  Die Trennung in Einzelstränge der doppelsträngigen DNA erfolgt nicht speziell in der Meiose I, sondern ist ein allgemeines Ereignis bei der DNA-Replikation und in der Zellteilung}
\end{todolist}
\subsection*{Eine Pflanze mit dem Genotyp AABbcc ist ...}
\begin{todolist} %\textit{\scriptsize }%
\item[\done] homozygot in Bezug auf zwei Loci.
\item heterozygot in Bezug auf zwei Loci. \textit{\scriptsize Pflanze ist heterozygot in Bezug auf einen Locus B (Bb)}
\item rezessiv in Bezug auf drei Loci. \textit{\scriptsize  'rezessiv' oder 'dominant' beziehen sich auf die Wirkung von Allelen in einem spezifischen phänotypischen Kontext und sind nicht direkt aus dem Genotyp ableitbar}
\item Alle drei Angaben stimmen. \textit{\scriptsize s.o.}
\end{todolist}
\subsection*{Manchmal kann ein Gen die Auspägung eines anderen Gens verhindern
oder überdecken. Diese genetische Wechselwirkung bezeichnet man als}
\begin{todolist}%\textit{\scriptsize }%
\item Genregulation.
\item Dominanz.
\item [\done]Epistasie.
\item zweite Paarregulation.
\end{todolist}
\subsection*{Vater und Mutter haben den Genotyp AaBbCcDdEE. Die Loci liegen alle
auf verschiedenen Chromosomen. Welcher Anteil der Nachkommen hat
(was die Loci A, B, C, D, und E angeht) denselben Genotyp wie die
Eltern?}
\[
P(\text{AA} \text{ und } \text{BB} \text{ und } \text{CC} \text{ und } \text{DD} \text{ und } \text{EE}) = \frac{1}{4} \times \frac{1}{4} \times \frac{1}{4} \times \frac{1}{4} \times 1 = \left(\frac{1}{4}\right)^4 \times 1 = \frac{1}{256}
\]
\begin{todolist}%\textit{\scriptsize }%
\item $75 \% (3/4)$
\item $25 \% (1/4)$
\item ca. $14 \% (9/64)$
\item $6,25 \% (1/16)$
\item[\done] ca. $0,4 \% (1/256)$
\item 0
\end{todolist}
\subsection*{Wenn alle Nachkommen einer Kreuzung heterozygot sind, sind die
Genotypen der Eltern wahrscheinlich ...}
\begin{todolist} %\textit{\scriptsize }%
\item [\done] AA, aa \textit{\scriptsize ergibt nur Aa }
\item Aa, Aa \textit{\scriptsize ergäbe AA (25$\%$), Aa (50$\%$) und aa (25$\%$)}
\item Aa, aa \textit{\scriptsize ergäbe Aa (50$\%$)
aa (50$\%$)}
\end{todolist}
\subsection*{Pleiotropie ist eine Genwirkung, bei der ...}
\begin{todolist}%\textit{\scriptsize }%
\item viele Gene auf ein Merkmal wirken. \textit{\scriptsize das wäre Polygenie}
\item mehrere Allele gleichzeitig exprimiert werden. \textit{\scriptsize  das wäre Codominanz oder Multiple Allelism}
\item beide Allele eines Locus die Ausprägung eines Merkmals
beeinflussen. \textit{\scriptsize beschreibt eher Allele Interaktion oder Kombination von Allelen in Bezug auf ein einzelnes Merkmal}
\item[\done] ein Gen auf viele Merkmale wirkt.
\item ein Allel die Expression des anderen Allels vollständig maskiert. \textit{\scriptsize das wäre Dominanz oder Epistase}
\end{todolist}
\subsection*{Welches der folgenden Paare kann KEIN Kind mit der Blutgruppe 0
haben? Angegeben sind die Blutgruppen der Eltern (Phänotypen).}
\begin{todolist} %\textit{\scriptsize }%
\item [\done] A x A
\item [\done] B x B
\item [\done] 0 x AB
\item Keines der drei Paare.
\end{todolist}
\subsection*{Eine Fruchtfliege, die heterozygot für graue Körperfarbe und normale
Flügel ist (phänotypisch Wildtyp), wird mit einer schwarzen Fliege mit
Stummelflügeln gekreuzt. Die Verteilung der Phänotypen unter den
Nachkommen ist: 778 Wildtyp, 785 schwarz/stummelflügelig, 162 grau/
stummelflügelig, 158 schwarz/normalflügelig. Wie gross ist der Abstand
der beiden betrachteten Gene in [cM] auf der Basis dieser Kreuzung?}
\[
   \text{Gesamtzahl der Nachkommen} = 778 + 785 + 162 + 158 = 1883
   \]
   \[
   \text{Rekombinanten Phänotypen} = 162 + 158 = 320
   \]
   \[
   \text{Rekombinierungsfrequenz} = \frac{320}{1883} \approx 0.170
   \]
   \[
   \text{Abstand in cM} = 0.170 \times 100 = 17 \text{ cM}
   \]
\begin{todolist} %\textit{\scriptsize }%
\item 0,17
\item 9
\item[\done] 17
\item 320
\item Zu wenig Informationen 
\end{todolist}
\subsection*{Thomas Bruder leidet an PKU, einer rezessiven Erbkrankheit. Die Eltern
des Bruders haben keine PKU. Wie wahrscheinlich ist es, das Thomas,
der in Bezug auf dieses Merkmal gesund ist, PKU übertragen kann?}
\[
\begin{array}{|c|c|}
\hline
\text{Elter 1} & \text{Elter 2} \\
\hline
\text{A} & \text{A} \\
\text{A} & \text{a} \\
\text{a} & \text{A} \\
\text{a} & \text{a} \\
\hline
\end{array}
\]
\[
\text{Genotypen: } \text{AA}, \text{Aa}, \text{aa}
\]
Die Wahrscheinlichkeiten sind:
\[
\text{P(AA)} = \frac{1}{4}
\]
\[
\text{P(Aa)} = \frac{1}{2}
\]
\[
\text{P(aa)} = \frac{1}{4}
\]
Da Thomas gesund ist, kann sein Genotyp entweder \( \text{AA} \) oder \( \text{Aa} \) sein. Die Wahrscheinlichkeit, dass Thomas ein Träger (\text{Aa}) ist, beträgt:
\[
\text{P(Träger)} = \text{P(Aa)} = \frac{1}{2} = 50\%
\]
\begin{todolist} %\textit{\scriptsize }%
\item $25 \% (1/4).$
\item $33,3 \% (1/3).$
\item [\done] $50 \% (1/2).$
\item $66,6 \% (2/3).$
\end{todolist}

\end{document}